% =============================================================================
% Ready-to-paste LaTeX: "The Relationship Between Musical Emotion and Film Genre"
% Uses citation keys already present in bib/library.bib. Placement: fits either as a
% subsection deepening Section~\ref{sec:fund_motivation_theory} (Fundamentals) or as the
% exploratory-analysis subsection in the Methods chapter (it reports the Cohen's d
% analysis that answers RQ2a). Adjust \subsection -> \section and the \label as needed.
% Cohen's d values are from the exploratory analysis on the 346-clip Set 1 corpus and are
% consistent with the Horror/Action figures already reported in the Fundamentals chapter.
%
% CITATION NOTE (added for the music-research strand, supervisor request #5:
% "Zusammenhang zwischen Emotion und Genre in der Literatur, auch in Musik").
% Four keys used below are NOT yet in bib/library.bib. Add:
%
%   @inproceedings{hu2007exploring,
%     title     = {Exploring Mood Metadata: Relationships with Genre, Artist and
%                  Usage Metadata},
%     author    = {Hu, Xiao and Downie, J. Stephen},
%     booktitle = {Proceedings of the 8th International Conference on Music
%                  Information Retrieval (ISMIR)},
%     pages     = {67--72}, year = {2007}
%   }
%
%   @inproceedings{laurier2009mood,
%     title     = {Music Mood Representations from Social Tags},
%     author    = {Laurier, Cyril and Sordo, Mohamed and Serr{\`a}, Joan and
%                  Herrera, Perfecto},
%     booktitle = {Proceedings of the 10th International Society for Music
%                  Information Retrieval Conference (ISMIR)},
%     pages     = {381--386}, year = {2009}
%   }
%
%   @article{eerola2011genrespecific,
%     title   = {Are the Emotions Expressed in Music Genre-specific? An Audio-based
%                Evaluation of Datasets Spanning Classical, Film, Pop and Mixed Genres},
%     author  = {Eerola, Tuomas},
%     journal = {Journal of New Music Research},
%     volume  = {40}, number = {4}, pages = {349--366}, year = {2011},
%     doi     = {10.1080/09298215.2011.602195}
%   }
%
%   @article{saari2016genre,
%     title   = {Genre-Adaptive Semantic Computing and Audio-Based Modelling for
%                Music Mood Annotation},
%     author  = {Saari, Pasi and Fazekas, Gy{\"o}rgy and Eerola, Tuomas and
%                Barthet, Mathieu and Lartillot, Olivier and Sandler, Mark},
%     journal = {IEEE Transactions on Affective Computing},
%     volume  = {7}, number = {2}, pages = {122--135}, year = {2016},
%     doi     = {10.1109/TAFFC.2015.2462841}
%   }
%
% NOTE: eerola2011genrespecific is a DIFFERENT paper from eerola2011comparison (the
% dataset paper, Eerola & Vuoskoski). Both are by the same first author and both from
% 2011; keep the two keys distinct.
% =============================================================================

\subsection{The Relationship Between Musical Emotion and Film Genre}
\label{subsec:emotion_genre}

The premise that film music encodes genre-appropriate emotion, established on
theoretical grounds in Section~\ref{sec:fund_motivation_theory}
\citep{gorbman1987unheard,juslin2013everyday}, raises a more specific empirical
question: \emph{which} emotional qualities characterise \emph{which} genres, and how
consistently? Both the prior literature and the exploratory analysis conducted for this
thesis converge on a clear, if uneven, answer.

\subsubsection*{Evidence from music research}

The link between expressed emotion and genre is not specific to film music; it is a
long-standing finding in music information retrieval, where it has been approached from
both directions. Working from metadata rather than audio,
\citet{hu2007exploring} analysed mood, genre, artist and usage descriptors drawn from
three independent sources (AllMusicGuide, epinions.com and Last.fm) and ran an
association test over 3{,}903 albums and 22 genres, yielding 7{,}134 genre--mood pairs.
They report substantial and systematic consistencies in the genre--mood relationship, and
it was precisely this structure that motivated their proposal to derive standardised
mood \emph{clusters} from the data -- the construction that subsequently became the
ground truth of the MIREX Audio Mood Classification task. \citet{laurier2009mood}
arrived at a compatible picture from social tags, deriving a semantic mood space from
Last.fm data by latent semantic analysis and showing that its principal structure
recovers the valence--arousal organisation of \citet{russell1980circumplex}, with mood
tags clustering in ways that align with musical style.

The converse direction -- whether emotion models are themselves genre-dependent -- was
tested directly by \citet{eerola2011genrespecific}, by the same author as the corpus used
in this thesis. Extracting 39 audio features from nine datasets spanning classical, film,
popular and mixed-genre music, he compared models trained and tested \emph{within} a
genre against models trained in one genre and applied to another. Performance within
genre reached 43\% of variance explained for valence and 62\% for arousal, but fell to
16\% and 43\% respectively across genres. Two conclusions follow, and both bear directly
on the design of this thesis. First, emotion and genre are so entangled that the mapping
from audio to emotion \emph{changes} with genre -- which is another way of saying that
emotion carries genre information, the premise of the present work. Second, that
entanglement sets a concrete expectation for cross-corpus transfer: an emotion model
carried to a new corpus should lose accuracy, and lose considerably more of it on valence
than on arousal. \citet{saari2016genre} close the loop by exploiting the dependency
constructively, showing that mood-annotation models made \emph{genre-adaptive} outperform
otherwise identical genre-agnostic ones on valence, arousal and tension.

Taken together, the music literature establishes the emotion--genre relationship as
bidirectional and reliable, but also as genre-conditional rather than universal. This is
what justifies treating emotion as an intermediate representation for genre prediction,
while cautioning that an emotion model estimated on one repertoire cannot be assumed to
behave identically on another -- the point examined empirically in
Chapter~
ef{ch:evaluation} through the cross-corpus transfer experiment.

\subsubsection*{Evidence from film music}

Empirical work on genre classification from film music has repeatedly observed that the
strength of the music--genre link is not uniform across genres. \citet{austin2010characterization}
report that high-intensity genres such as Action and Horror are markedly more separable
by their musical scores than narrative-driven genres such as Drama and Romance, and
attribute the residual confusion to the fact that individual cues within a labelled genre
may musically resemble a different genre. A parallel observation recurs in the music
emotion recognition (MER) literature: the surveys of \citet{kim2010music} and
\citet{kang2025there} document the persistent ``valence problem'', whereby arousal is
recovered from audio far more reliably than valence. This is directly relevant to genre,
because several genres are distinguished less by their level of arousal than by their
valence and by their discrete-emotion content. As \citet{eerola2011comparison} note for
the very corpus used here, Action and Horror occupy a nearly identical region of the
valence--arousal space -- negative valence combined with high arousal -- so a purely
dimensional description conflates them; it is the discrete emotion categories, and fear
in particular, that carry the information needed to separate them.

\subsubsection*{Evidence from the present corpus}

To quantify these relationships on the present corpus, the effect size (Cohen's $d$) of
each of the eight emotion ratings was computed between the clips belonging to a given
genre and those not belonging to it. The dominant emotional signature of each genre is
summarised in Table~\ref{tab:emotion_genre_signatures}. A coherent picture emerges. Fear
is the single strongest discriminator across the genre set: it sharply separates the
threat-oriented genres (Horror, $d\approx+0.75$; Crime, $d\approx+0.40$) from the
light-hearted ones (Comedy, $d\approx-0.61$). Valence forms a second axis, distinguishing
the more positively toned genres (Drama, $d\approx+0.57$; Comedy) from the negatively
toned ones (Horror, $d\approx-0.63$; Action). Each of these genres thus possesses an
interpretable affective fingerprint that is consistent with film-scoring convention:
Horror is defined by fear, Comedy by the \emph{absence} of fear together with positive
valence, Drama by elevated valence, and Action by high tension combined with an absence
of cheerfulness.

\begin{table}[t]
  \centering
  \caption{Dominant emotional signature of each genre on Set~1, measured by the effect
    size (Cohen's $d$) of the emotion rating between clips of that genre and all others.
    Positive $d$ indicates the emotion is elevated for the genre; negative $d$ indicates
    it is suppressed.}
  \label{tab:emotion_genre_signatures}
  \begin{tabular}{llr}
    \hline
    Genre       & Dominant emotion (direction) & Cohen's $d$ \\
    \hline
    Horror      & fear (high)          & $+0.75$ \\
    Biography   & happiness (high)     & $+0.67$ \\
    Comedy      & fear (low)           & $-0.61$ \\
    Drama       & valence (high)       & $+0.57$ \\
    Documentary & anger (low)          & $-0.51$ \\
    Action      & happiness (low)      & $-0.50$ \\
    Crime       & fear (high)          & $+0.40$ \\
    Adventure   & --- (none $|d|>0.25$) & $\phantom{+}0.25$ \\
    \hline
  \end{tabular}
\end{table}

The same analysis also delimits where the relationship breaks down. Adventure exhibits no
distinctive emotional signature -- no single emotion reaches an effect size of $|d|=0.25$
-- which reflects its tendency to co-occur with, and share the affective palette of,
several other genres rather than possessing one of its own. This absence is not a
measurement artefact but a substantive property of the genre, and it anticipates the
finding (Chapter~\ref{ch:evaluation}) that Adventure is the genre least predictable from
emotional features. More broadly, the uneven strength of the signatures in
Table~\ref{tab:emotion_genre_signatures} mirrors the asymmetry first reported by
\citet{austin2010characterization}: genres with a pronounced, conventionalised affective
identity (Horror, Comedy, Drama, Action) are recoverable from emotion, whereas genres
defined more by narrative or setting than by a characteristic emotional tone are not.

Taken together, these results support the use of perceived emotion as an
\emph{interpretable} intermediate representation rather than merely a predictive one. The
per-genre signatures are not only statistically detectable but also semantically
transparent and consistent with established film-music practice, so that a genre
prediction grounded in them can be explained in terms of the emotions the score expresses
-- an explanatory link that a direct mapping from low-level audio to genre does not
afford.
